\documentclass[border=3pt]{standalone}

% Circuits
\usepackage[european, straightvoltages, RPvoltages, americanresistor, americaninductors]{circuitikz}
\tikzset{every picture/.style={line width=0.3mm}}

% Notation
\usepackage{amsmath}

% Tikz Library
\usetikzlibrary{calc}

% Bipoles Specifications
\ctikzset{bipoles/thickness=1.5, label distance=1.2mm, voltage shift = 1.2}

% Arrows Above Components
\newcommand{\fixedvlen}[3][0.5cm]{% [semilength]{node}{label}
    \coordinate (#2-Vcenter) at ($(#2-Vfrom)!0.5!(#2-Vto)$);
    \draw[-{Stealth[scale=1.2]}, color=blue!85!black, ultra thick] ($(#2-Vcenter)!#1!(#2-Vfrom)$) -- ($(#2-Vcenter)!#1!(#2-Vto)$);
    \node[ anchor=\ctikzgetanchor{#2}{Vlab}, color=teal!90!black, font=\bfseries ] at (#2-Vlab) {#3};
}
\newcommand{\fixedvlendashed}[3][0.8cm]{% [semilength]{node}{label}
    \coordinate (#2-Vcenter) at ($(#2-Vfrom)!0.5!(#2-Vto)$);
    \draw[dash pattern=on 2pt off 2pt, -{Stealth[scale=1.1]}, color=red!80!black, very thick] ($(#2-Vcenter)!#1!(#2-Vfrom)$) -- ($(#2-Vcenter)!#1!(#2-Vto)$);
    \node[ anchor=\ctikzgetanchor{#2}{Vlab}, color=red!80!black, font=\bfseries ] at (#2-Vlab) {#3};
}

\begin{document}

\begin{circuitikz}
    % Coordinates
    \coordinate (GND) at (0,-1);
    \coordinate (N1) at (0,0);
    \coordinate (N2) at (0,3);
    \coordinate (N3) at (7,3);
    \coordinate (N4) at (7,0);

    % Circuit
    \draw
        (GND) node[ground]{} -- (N1)
        to[battery1, i=$I_Q$, v>, name=vbatt, color=teal!80!black, thick] (N2)
        to[R, l_=$R_A$, v^<, name=vrA, color=blue!80!black, ultra thick] (N3)
        to[R, l_=$R_B$, v^<, name=vrB, *-*, color=purple!85!black, thick] (N4)
        -- (N1);

    % Voltages
    \fixedvlen[0.5cm]{vbatt}{$V_{Q}$}
    \fixedvlen[0.7cm]{vrA}{$V_{A}$}
    \fixedvlen[0.7cm]{vrB}{$V_{B}$}
\end{circuitikz}

\end{document}